\documentclass[../../main.tex]{subfiles}

\begin{document}

\section{Obiettivi raggiunti} \label{sec:obiettivi-raggiunti}
  %    - propedeuticità al riempimento dei caf con roba vera
  %    - validazione input
  %    - testare analisi dati su oggetti completi
  \begin{figure}[ht]
    \centering
    \includegraphics[width=.99\textwidth]{fake-reco/fake-reco-v0}
    \caption{
      Schema di un \textit{workflow} del ufw che utilizza il modulo di fake-reco
      nella sua versione attuale.
    }
    \label{fig:fake-reco-v0}
  \end{figure}
  Al momento della stesura di questa tesi fake-reco, il cui utilizzo in un
  \textit{workflow} tipo è schematizzato in figura~\ref{fig:fake-reco-v0}, ha
  raggiunto i seguenti obiettivi / offre le seguenti funzionalità:

  \paragraph{Propedeuticità}
    I CAF di DUNE sono delle strutture dati nuove e ancora in via di sviluppo
    che i ricercatori devono imparare a usare, riempirli con la verità Monte
    Carlo permette di studiarli a fondo ed eventualmente di individuare
    storture nel loro design.

  \paragraph{Validazione dell'input}
    La maggior parte dei \textit{tool} di analisi dati sarà sviluppata in
    maniera naturale per lavorare sui CAF; avere un layer di traduzione
    uno a uno tra l'output di EDep-Sim e i CAF elimina la necessità di creare
    strumenti appositi per i dati simulati e, nel malaugurato caso in cui
    lo standard dei formati in ingresso dovessero cambiare, riduce lo
    sforzo reimplementativo alla sola fake-reco.

  \paragraph{Test degli strumenti di analisi}
    Per poter testare i \textit{workflow} di analisi dati è necessario avere i
    dati su cui lavorare.
    Delle strutture dati riempite con informazioni realistiche seppur troppo
    perfette permettono di testare più a fondo gli strumenti di
    analisi rispetto a strutture riempite con dati nulli o completamente
    casuali.

\end{document}